\section{Methodology}
\label{sec:methodology}

\paragraph{Random Forest.}
Random Forest is an ensemble method that averages the outputs of multiple decision trees trained on bootstrapped data subsets~\cite{breiman2001randomforest}. It is a strong baseline for Adult because the data contain mixed numeric and categorical attributes, missing values handled by imputation, and potentially non-linear threshold effects such as changes around age, capital gain, or working hours. Tree ensembles can model such interactions without assuming linear separability and remain comparatively stable when some one-hot encoded columns are sparse or weakly informative. In this project, the search space includes $n\_estimators \in \{200, 500\}$, $max\_depth \in \{\texttt{None}, 20, 40\}$, and $min\_samples\_leaf \in \{1, 5\}$. The best validation configuration was 500 trees with maximum depth 20 and minimum leaf size 1.

\paragraph{Multi-Layer Perceptron.}
The second method is a fully connected neural network trained with back-propagation~\cite{rumelhart1986backprop}. MLPs can learn smooth, high-order decision boundaries after categorical expansion and feature scaling, which makes them a reasonable contrast to tree ensembles even though tabular problems are not always their strongest setting. Here the MLP tests whether a compact neural model can recover competitive performance from the same information while using much less memory at inference time. The explored configurations use hidden layer sizes $(128, 64)$ and $(256, 128)$, with $\alpha \in \{10^{-4}, 10^{-3}\}$, learning rate $\in \{10^{-3}, 5\times10^{-4}\}$, early stopping enabled, and a maximum of 100 training epochs. The best validation configuration used hidden layers $(128,64)$, $\alpha = 0.001$, and learning rate $0.001$.

\paragraph{Why these two methods.}
The assessment requires two methods with substantially different modelling assumptions. Random Forest and MLP satisfy this requirement because one relies on many local decision rules aggregated across trees, while the other learns a shared parametric representation through gradient-based optimisation. This makes the pair suitable for a meaningful comparison of prediction quality, computational cost, and interpretability on the same tabular task. It also avoids a weaker comparison between two closely related variants from the same model family.
